% Final advisor message supplied by Martin on 2026-11-01.
% Include immediately before the student's AI-use declaration.
\chapter*{A témavezetők üzenete}
\addcontentsline{toc}{chapter}{A témavezetők üzenete}
\begingroup
\selectlanguage{magyar}

\noindent\textbf{Tisztelt Bíráló, kedves Kolléga!}

Az AI-használati nyilatkozat értelmezéséhez néhány megjegyzést szeretnénk fűzni, és ismertetjük azt a munkamódszert, amelynek használatára idén hallgatóinkat ösztönöztük. Egyúttal rögzítjük, milyen szakmai ellenőrzést végeztünk konzulensként, és mit vállalunk a bemutatott eredményekkel kapcsolatban.

Meggyőződésünk, hogy a generatív mesterséges intelligencia felelős kutatási alkalmazása túlmutathat az irodalomkutatás előkészítésén, a nyelvi ellenőrzésen és a szorosan felügyelt kódgeneráláson. A nagy nyelvi modellekre (LLM-ekre) épülő eszközök kutatási asszisztensként is használhatók: javasolhatnak kutatási irányokat és alternatív megoldásokat, segíthetik a prototípusok elkészítését, illetve megfelelően körülhatárolt részfeladatokat önállóan is végrehajthatnak. E munkamódszer felelős elsajátítását a korszerű mérnöki és kutatói felkészültség részének, a hallgatók erre való felkészítését pedig konzulensi feladatunknak tekintjük.

Az általunk adott útmutatás a felhasználást konkrét szakmai feltételekhez kötötte. A kutatási célokat és a kezdeti megoldási elképzeléseket konzulensi közreműködéssel alakítottuk ki, és szakmai megalapozottságukat ellenőriztük. Ezek mentén a hallgatótól a munka felelősségteljes kézben tartását vártuk el: a részfeladatok és követelmények pontos megfogalmazását, az alternatívák mérlegelését, valamint a döntések szakmai indoklását, az LLM-ek ismert hibamódjainak figyelembevételével. Az AI javaslatai vizsgálandó lehetőségeket jelentettek; elfogadásukat szakmai érvekhez és ellenőrzéshez kötöttük.

Különösen hasznosnak tartjuk az LLM-ekkel folytatott szokratikus dialógust. Kérdések, ellenpéldák és alternatívák módszeres vizsgálatával egy kezdeti ötlet vagy sejtés pontosítható, a rejtett feltevések felszínre hozhatók, és új megoldási lehetőségek azonosíthatók. Az AI ilyen közreműködése segítheti a kutatást; a kapott javaslatok tudományos értékét azonban a kidolgozás és az igazolás határozza meg.

A gyakorlatban a hallgató LLM segítségével kódot, modellt, kísérleti futtatókörnyezetet és adatelemző szkriptet is készíthetett. Vezetett módon a formalizmusok kialakításához is használhatott ilyen eszközt. Az ellenőrzés módját és mélységét a feladat szakmai jelentőségéhez és az esetleges tévedés következményeihez igazítottuk. A dolgozat eredményeit meghatározó állításoknál, számításoknál és megvalósításoknál az ellenőrzésnek az AI válaszától független szakmai alapra kellett támaszkodnia. Egy újabb AI-válasz vagy az eredeti válasz meggyőző megfogalmazása önmagában nem számított igazolásnak.

Konzulensként személyesen és érdemben ellenőriztük a dolgozat lényegi megoldási elképzeléseit, az eredményeket meghatározó kódot, valamint a matematikai állításokat és levezetéseket. A bemutatott megoldások és következtetések helyességéért, a dolgozatban rögzített feltételek és érvényességi tartomány mellett, konzulensként szakmai felelősséget vállalunk. Ez a vállalás az általunk elvégzett ellenőrzésen alapul.

Az AI használatának módjától és mértékétől függetlenül ugyanazt a szakmai mércét alkalmaztuk. Az azonosított hibákat javítani kellett, az eredményeket érintő bizonytalanságokat pedig fel kellett oldani vagy korlátként egyértelműen feltüntetni. Az AI tévedése nem szolgálhat indokként hibás állítás vagy megalapozatlan következtetés elfogadására; a végleges dolgozat szakmai tartalmára a fenti felelősségvállalásunk kiterjed.

A törzsszöveg tekintetében azt javasoltuk, hogy legalább annak gondolatmenetét és vázlatát a hallgató alakítsa ki, az LLM pedig az angol nyelvű megfogalmazásban és a nyelvi javításban kapjon jelentősebb szerepet. Elfogadhatónak tartottuk azt is, hogy a hallgató által felvázolt ábrából LLM segítségével készüljön a dolgozatba illeszthető változat. Ezekben az esetekben is ellenőrizni kellett, hogy a végleges szöveg vagy ábra pontosan fejezi-e ki a szakmai tartalmat.

Kérjük, hogy a nyilatkozatban feltüntetett, adott esetben \enquote{jelentős}-nek tekinthető AI-használatból önmagában ne következtessen a hallgató önálló szakmai hozzájárulásának hiányára. A generált szöveg vagy kód mennyisége önmagában nem megfelelő mérőszáma a tudományos teljesítménynek. Ennek megítéléséhez a hallgató által kidolgozott megoldásokat, meghozott szakmai döntéseket, elvégzett vizsgálatokat és igazolt eredményeket, valamint az eredménynyilatkozatban bemutatott munkamegosztást együttesen érdemes figyelembe venni.

A hallgatóknak adott útmutatásunkat konzulensként vállaljuk. Bízunk benne, hogy a munkamódszer és az elvégzett szakmai ellenőrzés bemutatása segíti az alábbi tételes nyilatkozat értelmezését és a dolgozat érdemi értékelését.

\medskip
\noindent Tisztelettel és üdvözlettel:

\noindent\textbf{Dr. Kocsis Imre és Farkas Martin}\\
konzulensek
\par
\endgroup
